\problemname{Kvicksand}
\noindent

Sandra har flyttat till Sandsvall för att studera. Efter att ha tillbringat några dagar i den
nya staden har Sandra märkt att det är kvicksand på flera platser. 
Om man åker till någon av dessa platser så fastnar man i sanden och kommer inte därifrån.
Sandra vill ha reda på exakt var kvicksanden finns så att hon kan undvika dessa platser.
Som tur är har Sandsvall, till skillnad från Göteborg, en välfungerande reseplaneringsapp.
Med hjälp av reseplaneraren har Sandra tagit reda på hur lång tid det tar att åka från 
hennes bostad till var och en av de andra platserna i staden.

Sandsvall kan representeras av ett rutnät med $N$ rader och $M$ kolumner. Sandras bostad
är i en rutorna i detta rutnät. Det tar en minut att åka mellan två närliggande rutor (två 
rutor är närliggande om de delar en sida). I vissa rutor finns det kvicksand, och då går 
det inte att åka därifrån alls.

Låt $(r, c)$ vara rutan på rad $r$ och kolumn $c$. Du får givet talet $d_{r,c}$ för varje 
ruta i rutnätet, det minsta möjliga antalet minuter det tar att åka från Sandras bostad 
till rutan. Din uppgift är att hantera $Q$ frågor av följande slag:

\begin{itemize}
    \item \verb|? r c|: Avgör om det finns kvicksand på ruta $(r, c)$. Du ska svara 
    \verb|ja|, \verb|nej|, \verb|kanske|, eller \verb|fel|.
    \item \verb|! r c d|: Ändra värdet på $d_{r, c}$ till $d$.
\end{itemize}

\section*{Input}

Den första raden innehåller tre heltal $N$, $M$ och $Q$ ($1 \leq N,M,Q \leq 3 \cdot 10^5$, $N\cdot M \leq 3 \cdot 10^5$).

Följande $N$ rader innehåller $M$ heltal var. Den $i$te av dessa rader innehåller heltalen
$d_{i,1}, d_{i,2}, d_{i,3}, \dots d_{i,M}$ ($-1 \leq d_{i,j} \leq 3 \cdot 10^5$).

Det kommer finnas exakt en ruta som uppfyller $d_{r, c} = 0$, detta är Sandras bostad. Om 
$d_{r, c} = -1$ så innebär det att det inte går att ta sig till ruta $(r, c)$ alls.

Följande $Q$ rader innehåller frågorna som ska behandlas, på formatet ovan. Talen $r, c$ som 
förekommer i frågorna uppfyller $1 \leq r \leq N$ och $1 \leq c \leq M$. Talen $d$ som förekommer 
uppfyller $-1 \leq d \leq 3 \cdot 10^5$, $d \neq 0$. Sandras bostad kommer alltid ha avståndet 
noll, och inga andra rutors avstånd kommer ändras till noll. 

\section*{Output}

Skriv ut svaret på varje fråga på formen \verb|? r c|.
Om det inte finns något sätt som kvicksanden kan vara fördelad på som är konsekvent 
med de givna talen $d$ ska du skriva ut \verb|fel|.
Annars, om det måste finnas kvicksand på ruta $(r, c)$, skriv ut \verb|ja|.
Om det garanterat inte finns kvicksand på ruta $(r, c)$, skriv ut \verb|nej|.
Om det både kan och inte kan finnas kvicksand på ruta $(r, c)$, skriv ut \verb|kanske|.

\section*{Poängsättning}
Din lösning kommer att testas på en mängd testfallsgrupper.
För att få poäng för en grupp så måste du klara alla testfall i gruppen.

\noindent
\begin{tabular}{| l | l | p{12cm} |}
  \hline
  \textbf{Grupp} & \textbf{Poäng} & \textbf{Gränser} \\ \hline
  $1$    & $17$        & $N = 1$. \\ \hline
  $2$    & $28$        & $N,M,Q \leq 100$. \\ \hline
  $3$    & $20$        & Svaret är aldrig \verb|fel|. \\ \hline
  $4$    & $35$        & Inga ytterligare begränsningar. \\ \hline
\end{tabular}

\section*{Förklaring av exempelfall}
